\section{Axis Ablations}
\label{sec:axis-ablations}

The axis ablations test whether the observed \method{} gains in the NLP experiments come from adaptive learning-rate modulation, adaptive gradient-noise control, or their combination. All configurations in this section use \adamw{} with warmup-linear scheduling as the base setup. \method{}-LR adapts only the learning-rate multiplier while keeping gradient noise fixed at zero, \method{}-Noise adapts only the gradient-noise magnitude while keeping the learning-rate multiplier fixed at one, and \method{}-Dual adapts both axes.

\input{tables/nlp_axis_ablation_agnews_noisy}
\input{tables/nlp_axis_ablation_sst2}

As shown in \Cref{tab:agnews-noise20-axis}, the noisy \agnews{} ablation clearly identifies adaptive gradient-noise control as the dominant useful axis. The warmup-linear baseline reaches $93.19\%$ peak accuracy and $90.79\%$ final accuracy, with a $2.40$ percentage-point drop. \method{}-LR improves final accuracy to $91.53\%$ (+$0.74$~pp) and reduces the drop to $1.19$ percentage points, but it does not improve peak accuracy over the baseline. In contrast, \method{}-Noise reaches $93.47\%$ peak accuracy and $93.45\%$ final accuracy, improving final accuracy by $2.66$~pp over the baseline and almost eliminating the best-to-final drop. \method{}-Dual gives the highest peak and final accuracy, $93.55\%$ and $93.51\%$, only $0.06$--$0.08$~pp above \method{}-Noise. The bulk of the noisy-\agnews{} gain therefore comes from adaptive gradient-noise control, while adaptive learning-rate modulation contributes only a small additional effect on its own.

The \sst{} ablation in \Cref{tab:sst2-axis} shows the same qualitative pattern, although with smaller margins. \method{}-LR reaches $90.55\%$ peak accuracy and $90.21\%$ final accuracy, only modestly improving over the warmup-linear baseline. \method{}-Noise gives the strongest result, with $91.12\%$ peak accuracy and $90.57\%$ final accuracy. \method{}-Dual reaches $90.80\%$ peak and $90.44\%$ final accuracy, improving over the baseline but not over the noise-only variant. Thus, on \sst{}, adding adaptive learning-rate modulation to adaptive gradient-noise control does not improve the result.

Taken together, the NLP ablations show that \method{} should not be interpreted as a monolithic method whose full dual-axis form is always best. In the tested Transformer fine-tuning settings, adaptive gradient-noise control is the strongest individual axis, while adaptive learning-rate modulation is weaker. This contrasts with the final \cifar{} experiments, where the successful configuration uses adaptive learning-rate modulation and gradient-noise injection was disabled after preliminary experiments found it destabilizing. The results therefore support the interpretation of \method{} as a framework for adaptive exploration control rather than a fixed optimizer recipe. Its effectiveness depends on exposing a control axis that is useful for the current training regime.